\begin{problem}{$N$거리 건너기}{standard input}{standard output}{1 second}{256 megabytes}

한결이는 오늘도 학교에 간다. 왜냐하면 한결이는 졸업과 동시에 입학을 하기 때문이다.

한결이를 놀리고 싶어하는 OO이는 한결이의 졸업식에서 입학을 축하하는 만행을 저질렀다. 하지만 만행은 그것으로 끝나지 않았다.

OO이는 한결이가 등교할 때 지나가는 횡단보도를 골라, 이를 특별한 마법으로 ``$N$거리"로 바꿔버렸다. ``$N$거리"는 정$N$각형 모양의 교차로로, $1, 2, \cdots, N$번 인도가 시계 방향으로 배치되어 있고, 인도들은 아래와 같이 $N$개의 횡단보도로 연결되어 있다.

\begin{itemize}
    \item $i (1 \le i \le N-1)$번 횡단보도는 $i$번 인도와 $i+1$번 인도를 연결한다.
    \item $N$번 횡단보도는 $N$번 인도와 $1$번 인도를 연결한다.
\end{itemize}

횡단보도에는 신호등이 하나씩 있다. 신호등은 아래 규칙에 따라 돌아가면서 $1$초씩 초록불이 켜진다.

\begin{itemize}
    \item 처음에는 $A_1$번 신호등의 초록불이 켜진다.
    \item $A_i$번 신호등의 초록불이 꺼짐과 동시에 $A_{i+1}$번 신호등의 초록불이 켜진다.
    \item 단, $i = N$ 일 때는 $A_N$번 신호등의 초록불이 꺼짐과 동시에 $A_1$번 신호등의 초록불이 켜진다.
\end{itemize}

한결이는 \textbf{$A_1$번 신호등의 초록불이 켜짐과 동시에} $1$번 인도에서 출발해 몇 개의 횡단보도를 건너서 $M$번 인도로 가려고 한다. 시계 방향과 반시계 방향 중 어떤 방향으로 가는 게 더 빠른지 구해보자.

단, 한결이는 걸음이 매우 빠르기 때문에 인도와 횡단보도 위를 걷는 시간은 무시해도 된다.

\InputFile
첫째 줄에는 횡단보도의 개수를 의미하는 $N$, 한결이가 가야 하는 인도의 번호 $M$이 공백으로 구분되어 주어진다.

둘째 줄에는 $N$거리 신호등의 초록불이 켜지는 순서를 의미하는 $N$개의 정수 $A_1, A_2, \cdots, A_N$이 공백으로 구분되어 주어진다.

\OutputFile
한결이가 $1$번 인도에서 $M$번 인도로 가기 위해 반시계 방향을 선택하는 것이 빠르다면 $CCW$, 시계 방향을 선택하는 것이 빠르다면 $CW$를 출력한다.

\Examples

\begin{example}
\exmpfile{example.01}{example.01.a}%
\exmpfile{example.02}{example.02.a}%
\end{example}

\Note
\begin{itemize}
    \item $3 \le N \le 100\,000$
    \item $2 \le M \le N$
    \item $1 \le A_i \le N (1 \le i \le N)$
    \item 수열 $A$의 원소는 서로 다르다. 즉, $i \ne j$ 이면 $A_i \ne A_j$이다.
    \item 입력으로 주어지는 수는 모두 정수이다.
\end{itemize}

TODO: 예제 설명 추가

\end{problem}

